\editionPage{mode-emploi}{Apprendre avec ce livre}{Comprendre, chercher, vérifier, recommencer}{Un parcours complet pour l'année de Seconde}
\editionEtape{Le cours donne du sens}
Les définitions, lois et méthodes précisent les conditions d'application. Les exemples montrent comment passer du texte à un calcul ; les schémas accompagnent les phénomènes, les montages et les modèles. Les activités expérimentales proposent une démarche à discuter avec l'enseignant.
\editionEtape{Douze exercices par chapitre}
Les exercices 1 à 4 installent les bases, les 5 à 8 consolident les méthodes et les 9 à 12 demandent davantage de raisonnement. Les points de difficulté sont indiqués dans les en-têtes. Chercher avant de lire la solution : une correction comprise se vérifie en refaisant l'exercice sans aide.
\editionEtape{Des figures dans les deux sens}
Quand un tracé est demandé, l'énoncé fournit les données et le support utiles. La version corrigée ajoute les vecteurs, rayons ou projections nécessaires. Les légendes distinguent les schémas de principe des dessins à l'échelle. Les unités et les échelles font partie de la réponse.
\editionEtape{Approfondir et s'évaluer}
Vingt défis mobilisent plusieurs idées de chaque chapitre. Des indications progressives précèdent leurs corrigés. Six évaluations de 90 minutes, notées sur 20, permettent de travailler sans le cours ; les barèmes valorisent les étapes de raisonnement. Ce format est pédagogique : il n'existe pas d'épreuve nationale unique de Seconde.
\editionRetenir{Une bonne réponse associe une idée physique, une relation justifiée, un calcul avec unités et une conclusion liée à la question.}
\clearpage
